\documentclass[hidelinks,onefignum,onetabnum,final]{siamart220329}  % for arxiv
%\documentclass[review,hidelinks,onefignum,onetabnum,final]{siamart220329}  % for submission

\usepackage{amsfonts,amssymb}
\usepackage{graphicx}
\DeclareGraphicsExtensions{.eps,.pdf,.png,.jpg}

% Used for creating new theorem and remark environments
\newsiamremark{remark}{Remark}
%definition already defined \newsiamremark{definition}{Definition}
\newsiamremark{hypothesis}{Hypothesis}
\crefname{hypothesis}{Hypothesis}{Hypotheses}
\newsiamthm{claim}{Claim}
\newsiamremark{example}{Example}

\usepackage{bm}

%\usepackage{amsopn}
%\usepackage{bm,bbm,empheq,verbatim,fancyvrb}
%\usepackage{booktabs,multirow,xspace}

\usepackage{tikz}
\usetikzlibrary{decorations.pathreplacing}
\usetikzlibrary{graphs,quotes}

\usepackage[kw]{pseudo}
\pseudoset{%
left-margin=6mm,%
topsep=5mm,%
label=\footnotesize\arabic*,%
idfont=\texttt,%
ctfont=\textsl,%
ct-left=\qquad\qquad(,%
ct-right=),%
}

\newcommand{\eps}{\epsilon}
\newcommand{\RR}{\mathbb{R}}

\newcommand{\grad}{\nabla}
\newcommand{\Div}{\nabla\cdot}

\newcommand{\bn}{\mathbf{n}}
\newcommand{\bu}{\mathbf{u}}
\newcommand{\bw}{\mathbf{w}}

\newcommand{\bU}{\mathbf{U}}

\newcommand{\bPhi}{\bm{\Phi}}
\newcommand{\bzero}{\bm{0}}

\newcommand{\dx}{\,dx}
\newcommand{\bds}{\bm{ds}}

\newcommand{\cK}{\mathcal{K}}
\newcommand{\cV}{\mathcal{V}}
\newcommand{\cW}{\mathcal{W}}
\newcommand{\cX}{\mathcal{X}}
\newcommand{\cY}{\mathcal{Y}}
\newcommand{\cZ}{\mathcal{Z}}

\newcommand{\pp}{\mathrm{p}}
\newcommand{\qq}{\mathrm{q}}

\newcommand{\ip}[2]{\left<#1,#2\right>}
\newcommand{\Bform}[2]{B\big(\left(#1\right),\left(#2\right)\big)}

\DeclareMathOperator{\diag}{diag}

\newcommand{\XX}{\ding{55}}

\DeclareMathOperator*{\argmin}{arg\,min}

\newfloat{pseudofloat}{t}{xyz}[section]
\floatname{pseudofloat}{Algorithm}

\usepackage{todonotes}
\newcommand{\elb}[1]{\todo[inline, color=blue!20]{EB: #1}}


% Sets running headers as well as PDF title and authors
%\headers{FIXME}{E. Bueler}

\title{Snow piles on conveyor belts: \\ a simplified glacier evolution model}

\author{Ed Bueler and Thomas Surowiec}

%\date{\today}
%\thanks{Department of Mathematics and Statistics, University of Alaska Fairbanks, USA (\email{elbueler@alaska.edu}).}

\graphicspath{{figs/}}

\begin{document}

\maketitle

\begin{abstract}
This is a very preliminary note!  It considers how to solve the main complementarity problem for the surface elevation of a glacier, in the case where we replace the actual surface value of the ice-flow velocity by a fixed velocity field, via a Skorohod-type formula.
\end{abstract}

%\begin{keywords}
%FIXME
%\end{keywords}
%\begin{MSCcodes}
%FIXME
%\end{MSCcodes}

Consider snow which falls at some rate $a(t,x,y)$, over a time interval $t\in [0,\infty)$ and for locations $(x,y)$ within a fixed map-plane domain $\Omega \subset \RR^2$.  More precisely, suppose $a$ gives the \emph{signed} instantaneous surface mass balance rate \cite{Cogleyetal2011} in ice-equivalent meters per second.  This means that $a>0$ when snow is falling ($a=$ accumulation rate) and $a<0$ when melt is occurring ($a=$ ablation rate).  While in actuality there is a time scale for the conversion of snow into maximal-density ice, in this note we suppose that process is instantaneous, so falling snow immediately becomes flowing ice in a glacier.  Also note that the time scale at which melt water runs off-glacier is sufficiently short that we also regard it as instantaneous.

Suppose that the land is at elevation $b(x,y)$, which is assumed time-independent.  The value of $b$ is the bedrock elevation underneath the glacier, or it is the land surface elevation if the land is not ice-covered.  Under these conditions, as snow piles-up, according to the values $a$, it will form a glacier, which flows under the action of gravity, according to the shape of the mountains given by $b$.

Let $s(t,x,y)$ be the surface elevation of the evolving glacier.  A \emph{surface kinematical equation} (or ``condition'' \cite{GreveBlatter2009}) for the glacier applies at the surface $s$,
\begin{equation}
\frac{\partial s}{\partial t} - \bu|_s \cdot \bn_s - a = 0,  \label{eq:ske}
\end{equation}
but this is only true at points \emph{on the ice}.  Over the whole domain $\Omega$, including both ice-covered and ice-free locations, the surface can never go below the bed $b$.  Thus the mathematical model for the evolution of $s$ is the following complementarity problem:\footnote{Complementarity problems are the strong forms corresponding to variational inequalities \cite{KinderlehrerStampacchia1980}.}
\begin{equation}
s - b \ge 0, \quad
\frac{\partial s}{\partial t} - \bu|_s \cdot \bn_s - a \ge 0, \quad
(s - b) \left(\frac{\partial s}{\partial t} - \bu|_s \cdot \bn_s - a\right) = 0  \label{eq:ncp}
\end{equation}

In the above equations $\bn_s = \left<-\partial s/\partial x,-\partial s/\partial y,1\right>$ is an upward surface-normal vector and $\bu|_s$ denotes the (upper) surface value of the ice flow velocity.  The surface kinematical equation \eqref{eq:ske} can be used to determine the vertical motion of the surface $\partial s/\partial t$ from the surface-normal component of the ice flow velocity $\bu|_s \cdot \bn_s$ and the surface mass balance $a$.

In fact, the 3D ice flow velocity $\bu$ solves a Stokes-type problem within the ice, below the surface, which is a very-viscous fluid \cite{GreveBlatter2009}.  Thus a complete glacier evolution model would couple \eqref{eq:ncp} with equations determining the 3D velocity and pressure fields through a Stokes solution at every instant.  For such a model see equations (1.1)--(1.6) in \cite{Bueler2025}; complementarity problem \eqref{eq:ncp} is system (1.3) therein.

However, a severely simplified model arises if the coupled-solution velocity is replaced by a pre-determined velocity field $\bu(t,x,y) \in \RR^3$ which is defined for every $t\in [0,\infty)$ and $(x,y)\in\Omega$.  In this note we consider such a model, in which the velocity field $\bu$ in the surface kinematical complementarity problem is given as \emph{data}:
\begin{equation}
s - b \ge 0, \quad
\frac{\partial s}{\partial t} - \bu \cdot \bn_s - a \ge 0, \quad
(s - b) \left(\frac{\partial s}{\partial t} - \bu \cdot \bn_s - a\right) = 0  \label{eq:pilencpearly}
\end{equation}
We will also suppose $s=b$ along $\partial \Omega$, that is, the glacier's thickness $H=s-b$ is forced to go to zero at the boundary of the computation domain.  In summary, $\Omega,a(t,x,y),b(x,y),\bu(t,x,y)$ are data for this problem, and $s(t,x,y)$ is the solution.

The interior condition of complementarity problem \eqref{eq:pilencpearly} is a linear advection PDE, that is, the surface kinematical equation \eqref{eq:ske} is now a linear advection.  Denote $\bu=\left<u,v,w\right>$ for the velocity components, and from now on denote partial derivatives with subscripts.  Let $c = w + a$; this is data.  Thus the interior PDE, which applies at ice-covered locations where $s>b$, is
\begin{equation}
s_t + u s_x + v s_y = c \label{eq:pilepde}
\end{equation}
Here $u(t,x,y)$, $v(t,x,y)$, and $c(t,x,y)$ are scalar functions given as data.  With this same notation, complementarity problem \eqref{eq:pilencpearly} is written
\begin{subequations}
\label{eq:pilencp}
\begin{align}
s - b &\ge 0 \\
s_t + u s_x + v s_y - c &\ge 0 \\
(s - b) \left(s_t + u s_x + v s_y - c\right) &= 0.
\end{align}
\end{subequations}

The interior condition \eqref{eq:pilepde} can be solved by characteristics \cite[section 3.2]{Evans2010}.  We state that result before modifying it to account for the inequalities in the complementarity problem.

\begin{lemma} \label{lem:chars}
Suppose $u, v, c \in C^1([0,T] \times \Omega)$ and let $s_0 \in C^1(\Omega)$ be initial data.  For each $(x_0,y_0) \in \Omega$, let
    $$t \mapsto (X(t),Y(t)) = (X(t;x_0,y_0), Y(t;x_0,y_0))$$
denote the solution of the characteristic ODE system
\begin{equation}
\frac{dX}{dt} = u(t,X,Y), \quad \frac{dY}{dt} = v(t,X,Y), \quad
X(0) = x_0, \quad Y(0) = y_0. \label{eq:chars}
\end{equation}
Then the unique classical solution of \eqref{eq:pilepde} satisfies
\begin{equation}
s\!\left(t, X(t), Y(t)\right) = s_0(x_0,y_0)
  + \int_0^t c\!\left(\tau, X(\tau), Y(\tau)\right) d\tau. \label{eq:charsol}
\end{equation}
\end{lemma}

\begin{proof}
Differentiate \eqref{eq:charsol} with respect to $t$ using the chain rule and
\eqref{eq:chars}, so that \eqref{eq:pilepde} is satisfied.
\end{proof}

If we ignored the inequalities of \eqref{eq:pilencp} then formula \eqref{eq:charsol} could become an explicit formula for $s(t,x,y)$ supposing that the characteristics could be inverted.  That is, we get a formula for $s(t,x,y)$ if $x=X(t;x_0,y_0)$ and $y=Y(t;x_0,y_0)$ can be solved for $x_0,y_0$ in terms of $t,x,y$.  In that case we could rewrite \eqref{eq:charsol} as
\begin{equation}
s(t,x,y) = s_0(x_0,y_0)
  + \int_0^t c\!\left(\tau, X(\tau;x_0,y_0), Y(\tau;x_0,y_0)\right) d\tau \label{eq:charsolagain}
\end{equation}
where it is understood that $x_0=x_0(t,x,y)$ and $y_0=y_0(t,x,y)$.

In practice the domain $\Omega$ is bounded, and the characteristic curves either be lead back to locations $(x_0,y_0) \in \Omega$ at time $t=0$, or they intersect with the boundary at a nonnegative time $t_0$: \, $(X(t_0),Y(t_0)) \in \partial\Omega$.  In the latter case some boundary value must supply the value for $s$ at that ``stopping'' time.  For simplicity we will assume zero thickness at the boundary:
\begin{equation}
s(t_0,X(t_0),Y(t_0)) = b(X(t_0),Y(t_0)) \quad \text{ if } (X(t_0),Y(t_0)) \in \partial\Omega \text{ for } t_0\ge 0.  \label{eq:charbc}
\end{equation}

FIXME now the Skorohod formula \cite[Lemma 6.14]{KaratzasShreve1991}, from queuing theory \cite[section 1.4.1]{Stewart2011}?, gives a solution of the complementarity problem \eqref{eq:pilencp}

\begin{lemma} \label{lem:skoro}  FIXME NEEDS CHECKING IN MORE DETAIL

Suppose the conditions of Lemma 1 hold, additionally $b \in C^1(\Omega)$, and
$s_0 \ge b$ on $\Omega$.  For $(x_0,y_0) \in \Omega$, let $\sigma(t) = \sigma(t;x_0,y_0)$
denote the unconstrained solution \eqref{eq:charsol} of \eqref{eq:pilepde}.  Define
\begin{equation}
\ell(t) = \max\!\left(0,\;
    \sup_{0 \le r \le t}\!\Big[b\!\left(X(r),Y(r)\right) - \sigma(r)\Big]\right),
\label{eq:skoro}
\end{equation}
so that $\ell$ is non-decreasing and $\ell(0) = 0$.  Then $S(t) = \sigma(t) + \ell(t)$ is the unique solution satisfying
\begin{equation}
S - b(X,Y) \ge 0, \qquad \dot{S} - c \ge 0, \qquad \big(S - b(X,Y)\big)\,\big(\dot{S} - c\big) = 0,
\label{eq:odencp}
\end{equation}
In other words,
\begin{equation}
s(t,X(t),Y(t)) = S(t) = \sigma(t) + \ell(t)
\end{equation}
solves \eqref{eq:pilencp} along characteristic \eqref{eq:chars}.
\end{lemma}

\begin{proof}
Apply the Skorohod reflection formula \cite[Lemma 6.14]{KaratzasShreve1991}
to $h(t) = \sigma(t) - b(X(t),Y(t))$, the unconstrained ice thickness along the
characteristic.  Since $s_0 \ge b$ we have $h(0) \ge 0$, and the Skorohod
lemma yields a unique pair $(H,\ell)$ with $H = h + \ell \ge 0$, $\ell$
non-decreasing, $\ell(0)=0$, and $H\,d\ell = 0$.  Setting
$S = H + b(X,Y) = \sigma + \ell$ and using $\dot{\sigma} = c$ from
\eqref{eq:charsol} gives \eqref{eq:odencp}.
\end{proof}

Note that by the chain rule and \eqref{eq:chars},
\begin{align*}
\dot S(t) &= s_t(t,X(t),Y(t)) + u(t,X(t),Y(t))\, s_x(t,X(t),Y(t)) \\
  &\hspace{29mm} + v(t,X(t),Y(t))\, s_y(t,X(t),Y(t)),
\end{align*}
so \eqref{eq:odencp} is the same as \eqref{eq:pilencp}.

\appendix

\section{A fixed-point argument for the coupled Skorohod-Stokes problem}
\label{app:fixedpoint}

Lemma~\ref{lem:skoro} solves complementarity problem~\eqref{eq:pilencp}
when the velocity $(u,v,w)$ is prescribed as data.  In the full glacier
evolution model the velocity is not data---it solves a Stokes sub-model
within the evolving ice domain, so the velocity and surface elevation
are coupled.  This appendix sketches how one might build a fixed-point
existence argument for the coupled problem using the Skorohod formula,
and identifies the key analytical difficulties.  The tone is deliberately
speculative; the goal is to record what would need to be proved.

\subsection{The fixed-point map}

Fix a time horizon $T>0$ and an exponent $r > 2$.  Define the admissible
set of surface elevations
\begin{equation}
\cK = \left\{s \in W^{1,r}(\Omega) : s \ge b \;\text{on}\; \Omega,\;
         s = b \;\text{on}\; \partial\Omega\right\}.
\label{eq:app:admissible}
\end{equation}
Since $\Omega \subset \RR^2$ and $r > 2$, Sobolev's embedding theorem gives
$W^{1,r}(\Omega) \hookrightarrow C(\bar\Omega)$, so each $s\in\cK$ is
continuous.  For $s\in\cK$ define the 3D ice domain
\begin{equation}
\Lambda(s) = \{(x,y,z) : b(x,y) < z < s(x,y)\} \subset \Omega\times\RR;
\label{eq:app:domain}
\end{equation}
continuity of $s$ and $b$ makes this a Lipschitz domain.

Let $\cU$ be a Banach space of surface velocity fields
$\bU = (u,v,w): [0,T]\times\Omega\to\RR^3$, to be specified.  Define
$T = T_2 \circ T_1$, the composition of two steps:
\begin{itemize}
\item $T_1:\cU \to C([0,T];\cK)$: given $\bU \in \cU$, apply
Lemmas~\ref{lem:chars} and~\ref{lem:skoro} to produce the surface
elevation $s = T_1 \bU$.
\item $T_2:C([0,T];\cK) \to \cU$: given $s\in C([0,T];\cK)$, for
each $t\in[0,T]$ solve the non-Newtonian Stokes sub-model
\cite{GreveBlatter2009} on the domain $\Lambda(s(t,\cdot))$,
and return the surface-trace velocity $\bu|_s(t,\cdot)$ extended by
zero to all of $\Omega$.
\end{itemize}
A fixed point $\bU = T\bU$ is a velocity field consistent with the
full coupled Stokes--surface kinematical complementarity problem.

\subsection{The Skorohod step \texorpdfstring{$T_1$}{T1}}

\paragraph{Input regularity needed.}
Lemma~\ref{lem:chars} requires $u,v \in C^1([0,T]\times\Omega)$ for
the characteristic ODE \eqref{eq:chars} to be classically well-posed.
A weaker sufficient condition is given by the DiPerna--Lions theory
of renormalized flows \cite{DiPernaLions1989}: a unique Lagrangian
flow exists for almost every initial point whenever
$(u,v) \in L^1(0,T;W^{1,1}(\Omega))$ and
$\mathrm{div}(u,v) \in L^1([0,T]\times\Omega)$.
Stokes incompressibility ($\Div\bu=0$ inside $\Lambda$) will be
relevant here.

\paragraph{Spatial regularity of $s$.}
The Jacobian of the characteristic flow map
$\Phi_t:(x_0,y_0)\mapsto(X(t;x_0,y_0),Y(t;x_0,y_0))$ satisfies the
Liouville identity
\begin{equation}
J(t) = \exp\!\left(\int_0^t \mathrm{div}(u,v)\big|_{\Phi_\tau} \, d\tau\right),
\label{eq:app:jacobian}
\end{equation}
so $J(t)>0$ when $\mathrm{div}(u,v) \in L^1$, ensuring invertibility.
When $(u,v)\in L^1(0,T;W^{1,r}(\Omega))$, regularity theory for ODE
flows gives $\Phi_t \in W^{1,r}(\Omega)$, and then the reassembled
surface elevation satisfies
\begin{equation}
s(t,\cdot) \in W^{1,r}(\Omega), \qquad
\|s(t,\cdot)\|_{W^{1,r}} \le
C\!\left(s_0,\, \|\bU\|_{L^1(0,T;\,W^{1,r})},\, \|a\|_{L^1(0,T;\,W^{1,r})}\right),
\label{eq:app:T1bound}
\end{equation}
provided $s_0\in W^{1,r}(\Omega)$ and $c = w+a \in L^1(0,T;W^{1,r}(\Omega))$.

\begin{remark}
The Skorohod correction $\ell(t;x_0,y_0)$ defined by \eqref{eq:skoro}
is a running maximum over $r\in[0,t]$, so its spatial regularity in
$(x_0,y_0)$ requires extra care near times when the maximum is
attained (i.e., where $S = b\circ\Phi$, the glacier margin).
A complete proof of \eqref{eq:app:T1bound} would need to verify,
for instance, that the contact set has measure-zero level sets,
or work with distributional derivatives.
\end{remark}

\subsection{The Stokes step \texorpdfstring{$T_2$}{T2}}

For $s\in\cK\subset W^{1,r}(\Omega)$ with $r>2$, the domain
$\Lambda(s)$ is Lipschitz and the regularized non-Newtonian Stokes
system is well-posed \cite{JouvetRappaz2011}.
Let $\pp = (n+1)/n$ denote the Glen flow exponent; for
$n=3$ one has $\pp=4/3$.  The solution velocity
$\bu\in W_0^{1,\pp}(\Lambda;\RR^3)$ satisfies an a priori bound
\begin{equation}
\|\bu\|_{W^{1,\pp}(\Lambda)} \le C_1,  \label{eq:app:apriori}
\end{equation}
where $C_1$ depends on physical constants and the geometry of
$\Lambda(s)$, in particular on $\|s\|_{W^{1,r}}$; see \cite{Bueler2025}.
A trace inequality then gives
\begin{equation}
\int_{\Gamma_s} |\bu|_s|^{\pp} \, dS \le C_1,  \label{eq:app:trace}
\end{equation}
and extended by zero, the surface trace satisfies
$\bu|_s(t,\cdot) \in L^{\pp}(\Omega;\RR^3)$ for each $t$.

\paragraph{What $T_2$ delivers.}
$T_2$ returns $\bu|_s \in L^{\pp}(\Omega;\RR^3)$, i.e.~velocity
components in $L^{4/3}(\Omega)$ for $n=3$.
Crucially, $T_2$ does \emph{not} directly deliver $W^{1,q}(\Omega)$
regularity for any $q > 0$.

\subsection{The regularity gap}

For the cycle
$\cU \xrightarrow{T_1} C([0,T];\cK) \xrightarrow{T_2} \cU$
to close, $\cU$ must satisfy both the input requirement of $T_1$
and the output property of $T_2$.  Since $T_1$ requires at minimum
$(u,v)\in L^1(0,T;W^{1,1}(\Omega))$ (for DiPerna--Lions) but
$T_2$ delivers only $\bu|_s \in L^1(0,T;L^{\pp}(\Omega))$
with $\pp=4/3$, the cycle does not close:
\[
W^{1,r}\text{ velocity} \;\xrightarrow{T_1}\; W^{1,r}\text{ surface}
\;\xrightarrow{T_2}\; L^{\pp}\text{ velocity}
\;\not\hookrightarrow\; W^{1,1}\text{ velocity}.
\]

There is also an exponent tension.  The surface elevation space requires
$r > 2$ (for $\Lambda(s)$ to be a Lipschitz domain via Sobolev
embedding), while the Stokes law has exponent $\pp = 4/3 < 2$.
Even if $\bu|_s \in W^{1,\pp}(\Omega)$ could be established, one has
$W^{1,\pp} \not\subset W^{1,r}$ for $r > \pp$, so the input
requirement of $T_1$ would still not be met by the output of $T_2$.

\subsection{Hypotheses that would close the argument}

Gaining spatial regularity of $\bu|_s$ beyond $L^{\pp}$ would
require higher interior regularity of the Stokes velocity:
$W^{2,\pp}$ regularity of $\bu$ inside $\Lambda(s)$ would give
(by a trace theorem) $\bu|_s \in W^{1,\pp}(\Omega)$.
Such $W^{2,\pp}$ regularity requires $\partial\Lambda(s)$ to be
$C^{1,\alpha}$ \cite{JouvetRappaz2011}, i.e.~$s\in C^{1,\alpha}(\Omega)$,
which in turn follows from $s\in W^{2,r}(\Omega)$ via Sobolev.
So higher surface elevation regularity would bootstrap Stokes
regularity, but the argument is circular without an independent
source of that higher regularity.

\begin{hypothesis}[Stokes surface trace regularity]
\label{hyp:stokestraceW1}
There exist exponents $r>2$ and $q>1$ such that for every
$s\in\cK \subset W^{1,r}(\Omega)$ the Stokes velocity surface trace
satisfies $\bu|_s \in W^{1,q}(\Omega;\RR^3)$, with
\begin{equation}
\|\bu|_s\|_{W^{1,q}(\Omega)} \le C\bigl(\|s\|_{W^{1,r}(\Omega)}\bigr).
\label{eq:app:hyp1}
\end{equation}
\end{hypothesis}

\noindent
For \emph{Newtonian} ice ($n=1$, $\pp=2$) on a $C^{1,1}$ domain,
classical $H^2$ regularity of the linear Stokes system gives
$\bu|_s \in H^1(\Omega)$, so Hypothesis~\ref{hyp:stokestraceW1} holds
with $q=2$ in that case.
For shear-thinning ice ($\pp<2$) on a domain that is merely Lipschitz,
this appears to be open.
A related domain-dependence statement for the Stokes map is
Conjecture~1 in \cite{Bueler2025}; that conjecture concerns
Lipschitz continuity rather than regularity, but both ultimately
rest on the same missing theory of how the Stokes solution varies
with the domain.

If Hypothesis~\ref{hyp:stokestraceW1} holds with $q\ge 1$, and if
the surface trace $(u,v) = (\bu|_s)_{1,2}$ satisfies
$\mathrm{div}(u,v) \in L^1$ (which follows from Stokes
incompressibility via the relation $w|_s = -\int_b^s \Div(u,v)\,dz$,
up to a boundary term at the margin), then the DiPerna--Lions
theory \cite{DiPernaLions1989} applies to the characteristic ODE
\eqref{eq:chars}, closing the regularity loop.

\begin{hypothesis}[Lipschitz continuity of Stokes trace in $s$]
\label{hyp:lipschitz}
Under the conditions of Hypothesis~\ref{hyp:stokestraceW1}, for every
$R>0$ there exists $C(R)>0$ such that for $s_1,s_2\in B_R\cap\cK$,
\begin{equation}
\|\bu|_{s_1} - \bu|_{s_2}\|_{W^{1,q}(\Omega)} \le
C(R)\,\|s_1-s_2\|_{W^{1,r}(\Omega)}.  \label{eq:app:hyp2}
\end{equation}
\end{hypothesis}

\noindent
Hypothesis~\ref{hyp:lipschitz} with $W^{1,q}$ replaced by $L^{r'}$
(where $1/r+1/r'=1$) is exactly Conjecture~1 in \cite{Bueler2025},
which is also the key ingredient for well-posedness of the implicit
time-step variational inequality in that reference.

\paragraph{Invariant set.}
For either the Schauder or Banach fixed-point theorem one needs a
bounded closed convex set $K\subset\cU$ with $T(K)\subseteq K$.
The a priori Stokes bound \eqref{eq:app:apriori} gives
$C_1=C_1(R_s)$ with $R_s = \|s\|_{W^{1,r}}$, and
\eqref{eq:app:T1bound} controls $R_s$ in terms of
$R_\cU=\|\bU\|_\cU$.
A self-consistent radius $R>0$ therefore exists if
\begin{equation}
C_1\bigl(C_{\mathrm{Sko}}(R,s_0,a)\bigr) \le R,
\label{eq:app:selfconsistent}
\end{equation}
where $C_{\mathrm{Sko}}$ is the bound from \eqref{eq:app:T1bound}.
For small time horizon $T$, or small forcing $a$, such a radius
can always be found.

\paragraph{Compactness vs.\ contraction.}
Under Hypotheses~\ref{hyp:stokestraceW1} and~\ref{hyp:lipschitz},
$T$ is continuous on bounded sets.
A \emph{Schauder fixed point} additionally requires $T$ to be compact;
compactness of $T_2$ follows from the uniform $W^{1,\pp}(\Lambda)$
bound on $\bu$ together with the compact Sobolev embedding
$W^{1,\pp}\hookrightarrow\hookrightarrow L^m$ for $m<\pp^*$,
but propagating that compactness through $T_1$ requires further
argument.
A \emph{Banach contraction mapping} requires
$\|T\bU_1-T\bU_2\|_\cU \le L\|\bU_1-\bU_2\|_\cU$ with $L<1$;
this would follow from Hypothesis~\ref{hyp:lipschitz} combined with
Gronwall-type ODE estimates for the Lipschitz continuity of $T_1$,
provided $T$ (or the time-step $\Delta t$) is small enough.

\subsection{Relation to the fully implicit approach in \cite{Bueler2025}}  \label{app:implicit}

The framework in \cite{Bueler2025} avoids the above regularity cycle
by formulating an \emph{implicit} time-step variational inequality (VI):
find the new surface elevation $s\in\cK$ such that the Stokes velocity
at $s$ and $s$ itself simultaneously satisfy the VI.
Well-posedness of this implicit step is established in \cite{Bueler2025}
(conditionally on two conjectures analogous to
Hypotheses~\ref{hyp:stokestraceW1}--\ref{hyp:lipschitz}) using the
theory of coercive monotone operators \cite{KinderlehrerStampacchia1980},
without any Picard iteration between Skorohod and Stokes steps.

The Skorohod-based iteration proposed here---alternating between
the Skorohod formula (with given velocity) and the Stokes solve
(with given surface)---is an alternative route to the same coupled
solution.
Convergence of the iteration follows when $T$ is a contraction; under
uniqueness of the implicit VI solution
(which follows from strict monotonicity under Conjecture~2 of
\cite{Bueler2025}), the fixed-point limit and the implicit VI solution
then coincide.

In summary, the central open problem for either approach is
Hypothesis~\ref{hyp:stokestraceW1}: that the Stokes velocity surface
trace on a Lipschitz glacier domain has $W^{1,q}$ regularity as a
function of the map-plane coordinates, with a bound that is Lipschitz
in the surface elevation.

\bibliographystyle{siamplain}
\bibliography{pile}

\end{document}
